\chapter{Experiments}
\label{ch:experiments}

\section{Methodology}
\label{sec:methodology}

We benchmark three algorithms -- brute-force NSP, Yen-NSP, and our
implementation of the Chen--Wein--Zhang (CWZ) algorithm -- on graphs
drawn from five generator families:
\begin{itemize}
  \item Erd\H{o}s--R\'enyi random digraphs $G_{n,p}$ with $p=0.30$;
  \item random DAGs (vertices in topological order; each forward edge
    present with probability $p=0.30$);
  \item layered digraphs (two endpoint vertices $s$, $t$ plus $L$ interior
    layers of width $W$, with forward-edge probability $p_f=0.40$ and
    back-edge probability $p_b=0.10$). The name refers to this generator's
    layer topology and \emph{not} to the $(s,t)$-layered property of
    Definition~4.1: these instances do not satisfy it -- a typical one has
    most of its vertices off every shortest $s\to t$ path and most of its
    back-edges not strictly backward -- so both reductions run in full on
    them, exactly as on the other families;
  \item grid digraphs (rows-by-columns lattice with rightward and downward
    edges only);
  \item \emph{diamond-chain} digraphs, a deliberately adversarial family
    described in Section~\ref{sec:adversarial}.
\end{itemize}

For each family we sweep $|V|\in\{6,8,10,12,16,20,25,30,40,50,60,70\}$, drawing
$20$ instances per size with deterministic seeds. Any instance in which $t$
is unreachable from $s$ is redrawn, so every measurement is taken on a graph
that actually poses the problem (Section~\ref{sec:existence}). All three
algorithms run across the whole sweep. The brute-force oracle is a branch-and-bound search
rather than an exhaustive enumeration: a reverse Dijkstra supplies the
admissible lower bound $d(u,t)$, and any prefix with
$\text{cost}+d(u,t)\ge\text{best}$ is pruned. That bound is what makes an
exponential-in-the-worst-case oracle usable across the sweep; it collapses,
by construction, only on graphs with very many equal-cost shortest paths,
where every shortest path must still be walked. Each algorithm call is
wrapped in a $30$-second wall-clock
budget; calls that exceed it are recorded as ``budget exceeded'' and
excluded from the timing aggregates. Edge weights are integers drawn
uniformly from $[1,10]$. All three algorithms are implemented in C++20,
compiled at \texttt{-O3}, and timed with wall-clock measurements taken
\emph{around the algorithm call only} (graph generation excluded). The
benchmark binary (\texttt{build/bench}) writes one CSV row per
$\langle\text{algorithm},\text{family},|V|,\text{trial}\rangle$ and the
plots are produced by \texttt{bench/plots/make\_plots.py}.

For each $(\text{family},|V|)$ cell we report the \emph{median} runtime
across the $20$ trials. Plots additionally show the inter-quartile range
($25^{\text{th}}$--$75^{\text{th}}$ percentile) as a shaded band around
the median, which is robust to single-trial outliers and to the
heavy-tailed runtimes that appear on dense Erd\H{o}s--R\'enyi instances.

\section{Runtime by family and size}
\label{sec:runtime-by-family}

Figure~\ref{fig:runtime} shows the runtime breakdown by family, with one
sub-panel per family. The vertical axis is logarithmic and the shaded
band around each line is the IQR. Three observations stand out.

\paragraph{All three algorithms are fast on the benign families.}
Even at $|V|=50$, CWZ on the densest family (Erd\H{o}s--R\'enyi) has a
median of $208$~ms at $|V|=70$ ($59$~ms at $|V|=50$). Brute force, despite its exponential worst case, is
the \emph{fastest} of the three on the four benign families: its lower-bound
pruning restricts the search to paths that could still beat the best
not-shortest path found so far, and on these inputs that is a handful of
short prefixes. Yen is likewise fast throughout, finishing in tens of
microseconds at $|V|=50$.

\paragraph{CWZ's runtime is family-sensitive.} On Erd\H{o}s--R\'enyi
graphs the median grows from $51$~\textmu s at $|V|=10$ to $208$~ms at
$|V|=70$, a four-orders-of-magnitude climb driven by the number of
back-edges that the inner enumeration must consider. On the random-DAG
family the same sweep climbs only from $11$~\textmu s to $3.4$~ms because
DAGs lack the dense back-edge structure of $G_{n,p}$. Layered and grid
graphs sit in between.

\paragraph{Yen dominates on the four random families.} For the input
distributions we tested, Yen-NSP terminates orders of magnitude faster
than CWZ. This is the expected behaviour on graphs where the number of
distinct simple shortest paths is small: Yen only has to step past a
handful of paths before finding one with strictly greater weight, while
CWZ pays its polynomial price uniformly. The interest of CWZ is precisely
in \emph{worst-case} polynomiality -- it is the only algorithm with a
fixed polynomial bound regardless of how many shortest paths the graph
contains. The next subsection shows that worst-case behaviour
empirically.

\paragraph{The diamond-chain panel inverts the picture.} The rightmost
sub-panel of Figure~\ref{fig:runtime} shows the same axes for the
diamond-chain family. Yen's curve climbs steeply -- from microseconds
at $|V|=10$ to $3.7$~s at $|V|=70$ -- while the CWZ curve stays
microsecond-flat at $51$~\textmu s. This is the empirical face of the
$2^k$ blow-up Yen must endure on graphs with exponentially many shortest
paths, against which CWZ's polynomial bound is the only protection.

\begin{figure}
  \centering
  \includegraphics[width=\textwidth]{../results/runtime_by_family.pdf}
  \caption{Median runtime versus $|V|$ across families (log $y$-axis), with
    inter-quartile range shaded. Each family is a sub-panel; the
    diamond-chain panel inverts the brute-vs-Yen-vs-CWZ ordering visible in
    the other four, with both exponential algorithms climbing past CWZ while
    CWZ stays flat. Twenty trials per $(\text{family},|V|)$ cell. Points
    sitting exactly on the $1$~\textmu s line are below our wall-clock timer
    resolution and are clamped there for display, so a flat segment at that
    level means ``too fast to measure'', not constant time.}
  \label{fig:runtime}
\end{figure}

\section{CWZ scaling}

Figure~\ref{fig:scaling} shows the CWZ median runtime alone on a log-log
scale, with all five families overlaid. Fitting a power law to the median
runtime of each family separately gives exponents that vary widely with
graph structure:

\begin{center}
\begin{tabular}{lc}
\toprule
family & log-log slope \\
\midrule
diamond chain  & $1.29$ \\
grid           & $2.28$ \\
layered        & $2.90$ \\
random DAG     & $2.90$ \\
Erd\H{o}s--R\'enyi & $4.34$ \\
\bottomrule
\end{tabular}
\end{center}

The spread tracks edge density: the sparse structured families sit near
$|V|^{1}$--$|V|^{3}$, while dense Erd\H{o}s--R\'enyi graphs
($|E|=\Theta(|V|^{2})$ at $p=0.30$) reach $\approx |V|^{4.3}$. The fits use each instance's
\emph{actual} vertex count, not the sweep target, which differ by up to
$30\%$ for the grid and diamond-chain families. Every
family remains far below the paper's worst-case bound -- for
$|E|=\Theta(|V|)$ the $|V|^{4}|E|^{3}$ bound is a slope of $\approx 7$ in
$\log|V|$, and for the dense Erd\H{o}s--R\'enyi regime it is
$\approx 10$. This indicates that the inner enumeration of our
implementation prunes aggressively in practice on these inputs, and that
the observed exponent is a property of the input family rather than a
single characteristic number.

\begin{figure}
  \centering
  \includegraphics[width=0.85\textwidth]{../results/cwz_scaling.pdf}
  \caption{CWZ median runtime on a log-log scale, all five families.}
  \label{fig:scaling}
\end{figure}

\section{Agreement between algorithms}
\label{sec:agreement}

The strongest check available is against the exact oracle. Because the
brute-force search now runs across the entire sweep, every one of the
$1200$ benchmark instances can be compared against it directly:
\textbf{CWZ agrees with brute force on all $1200$}, spanning all five
families and $|V|$ up to $81$. This extends the correctness evidence well
beyond the $|V|\le 14$ reach of the fuzz harness of
Chapter~\ref{ch:implementation}.

Against Yen the picture is nearly as clean. The script
\texttt{bench/plots/make\_plots.py} emits
\texttt{results/agreement\_summary.csv} which records, per
$(\text{family},|V|)$ cell, how many of the $20$ trials had matching NSP
costs. Agreement is $100\%$ on the four benign
families across the whole sweep, and on the diamond-chain family up to
$|V|=30$.

The only cells with less than full agreement are the two largest
diamond-chain sizes, $|V|=40$ and $|V|=50$ -- and there it is
\emph{Yen}, not CWZ, that is wrong. On those instances the number of
distinct shortest paths ($2^{k}$ with $k\ge 13$) exceeds Yen's
enumeration cap $k_{\max}=5000$, so Yen exhausts its budget and reports
``no NSP'' even though one exists. CWZ returns the correct NSP (cost
$2k+1$) in tens of microseconds on the same instances. This is a sharper
form of the adversarial story of Section~\ref{sec:adversarial}: a
capped Yen does not merely run slowly on graphs with exponentially many
shortest paths, it silently returns the wrong answer.

This agreement check is complementary to the fuzz harness of
Chapter~\ref{ch:implementation}: not only does our CWZ implementation
match the brute-force oracle on small inputs, it matches Yen on every
benign instance and \emph{beats} Yen (matching brute force) on the
adversarial instances where Yen's cap fails.

\begin{table}
  \centering
  \begin{tabular}{lrrr}
  \toprule
  family & $|V|$ range & trials per cell & agree \\
  \midrule
  Erd\H{o}s--R\'enyi & 6--50 & 20 / size & 100\% \\
  random DAG        & 6--50 & 20 / size & 100\% \\
  layered           & 6--50 & 20 / size & 100\% \\
  grid              & 6--50 & 20 / size & 100\% \\
  diamond chain     & 6--30 & 20 / size & 100\% \\
  diamond chain     & 40--50 & 20 / size & Yen capped$^{*}$ \\
  \bottomrule
  \end{tabular}
  \caption{CWZ vs Yen agreement on benchmark instances. The summary CSV
    \texttt{results/agreement\_summary.csv} carries the per-cell
    breakdown. $^{*}$On diamond chains at $|V|\ge 40$ Yen exceeds its
    $k_{\max}=5000$ enumeration cap and reports no NSP; CWZ returns the
    correct NSP, matching brute force, in tens of microseconds.}
  \label{tab:agreement}
\end{table}

\section{NSP existence}
\label{sec:existence}

A practical caveat in interpreting Figure~\ref{fig:runtime}: on some of the
smallest generated instances no NSP exists at all. Figure~\ref{fig:existence}
reports the fraction on which one does, as determined by the brute-force
oracle.

Two effects were at work here, and only one of them is interesting. The first
was an artefact: a graph in which $t$ is simply unreachable from $s$ has no
shortest path, so the question of a next-to-shortest one is vacuous. Sparse
families produced such graphs often at small $|V|$ -- $13/20$
Erd\H{o}s--R\'enyi instances at $|V|=6$ -- and in the layered family the cause
was structural rather than statistical: with $\text{width}=2$ and
$p_f=0.40$, a pair of consecutive layers receives no forward edge at all with
probability $(1-p_f)^{\text{width}^2}=0.6^{4}\approx 13\%$, severing the
graph. The generator now guarantees at least one forward edge between every
consecutive pair of layers, and the benchmark redraws any instance in which
$t$ is unreachable from $s$. Every one of the $1200$ instances reported here
is therefore genuinely solvable.

What remains is real. Of those $1200$ instances, $1122$ have an NSP and $78$ do
not, every one of the latter at $|V|\le 16$. These are graphs with exactly one
simple $s\to t$ path -- which is then the shortest path by default, leaving
nothing to be next-to-shortest -- or, in a handful of cases, several paths
whose costs all tie at the minimum, so that none of them is
\emph{not}-shortest. Both are legitimate NSP instances on which the correct
answer is ``none'', and all three algorithms are required to say so; we keep
them deliberately, as they are the only coverage the benchmark has of that
code path. The effect disappears as soon as the graphs are large enough to
admit several distinct-cost routes: every family is at $100\%$ by $|V|=20$.

It is worth noting how this figure was obtained. An earlier version plotted
the fraction on which \emph{Yen} reported an NSP, which made the curve for
the diamond-chain family collapse to zero at $|V|\ge 40$ -- on the one
family that provably always has an NSP. What it was really measuring was
Yen exhausting its $k_{\max}$ enumeration cap
(Section~\ref{sec:agreement}). Existence is a property of the instance,
not of the solver, so it must be read off an exact oracle.

\begin{figure}
  \centering
  \includegraphics[width=0.75\textwidth]{../results/nsp_existence.pdf}
  \caption{Fraction of generated instances on which an NSP exists, as
    determined by the brute-force oracle.}
  \label{fig:existence}
\end{figure}

\section{Adversarial inputs: where CWZ wins}
\label{sec:adversarial}

The diamond-chain family in Section~\ref{sec:methodology} (which already appears as the
right-most panel of Figure~\ref{fig:runtime}) is repeated here with a
dedicated sweep over $k$ to make the $2^{k}$ scaling visible at higher
resolution.

\paragraph{Diamond-chain construction.} For a parameter $k$ we build a
layered digraph as follows. Vertex $0$ is $s$. For each $i\in\{1,\dots,k\}$
we add three vertices: $\text{top}_i$, $\text{bot}_i$, and $\text{merge}_i$;
$\text{merge}_k$ is $t$. Set $\text{merge}_0\equiv s$. For each $i$, add
four unit-weight edges:
$\text{merge}_{i-1}\to\text{top}_i$, $\text{merge}_{i-1}\to\text{bot}_i$,
$\text{top}_i\to\text{merge}_i$, and $\text{bot}_i\to\text{merge}_i$.
Finally add a direct edge $s\to t$ of weight $2k+1$.

The construction has $1+3k$ vertices, $4k+1$ edges, and exactly $2^{k}$
distinct simple shortest $s\to t$ paths (each of weight $2k$). The unique
NSP is the direct $s\to t$ shortcut of weight $2k+1$.

\paragraph{Yen's exponential blowup.} On this family Yen must enumerate
all $2^{k}$ shortest simple paths before reaching the NSP.
Table~\ref{tab:adv} shows the measured runtimes from the dedicated
\texttt{build/bench\_adversarial} sweep; Figure~\ref{fig:adversarial}
plots them.

\begin{table}
  \centering
  \begin{tabular}{rrrrr}
  \toprule
  $k$ & $|V|$ & $|E|$ & Yen (s) & CWZ (s) \\
  \midrule
  10 & 31 & 41 & 0.023 & 0.000014 \\
  12 & 37 & 49 & 0.332 & 0.000016 \\
  14 & 43 & 57 & 10.04 & 0.000018 \\
  15 & 46 & 61 & 56.18 & 0.000022 \\
  20 & 61 & 81 & --   & 0.000025 \\
  24 & 73 & 97 & --   & 0.000031 \\
  \bottomrule
  \end{tabular}
  \caption{Diamond-chain adversarial benchmark. Yen-NSP must enumerate
    all $2^{k}$ shortest paths and was halted after exceeding the
    $30$~s budget at $k=15$ (recorded budget overflow:
    $56$~s). CWZ stays in the tens of microseconds across the full
    sweep. The two algorithms are timed in separate passes and each CWZ
    figure is the median of $11$ repetitions; see the measurement note
    below.}
  \label{tab:adv}
\end{table}

\begin{figure}
  \centering
  \includegraphics[width=0.8\textwidth]{../results/adversarial.pdf}
  \caption{Adversarial diamond-chain: Yen-NSP's runtime grows by roughly a
    factor of five per increment of $k$ while CWZ stays microsecond-fast.
    The number of shortest paths Yen must enumerate doubles with $k$, but its
    spur loop rescans the whole result list for every spur index, so the
    measured growth is closer to $4^{k}$ than $2^{k}$.}
  \label{fig:adversarial}
\end{figure}

Between $k=10$ and $k=15$, Yen's runtime grows from $0.023$~s to
$56$~s -- consistent with the $2^{k}$ scaling we expected: each
increment of $k$ doubles the number of shortest paths Yen has to
enumerate before finding a strictly longer one. CWZ, in contrast, finds
the NSP in microseconds across the full sweep because the graph contains
exactly one back-edge (the $s\to t$ shortcut). That edge goes forward in
distance, so \textsc{NextSP-Straight} removes it and records it as a
single candidate (Section~\ref{sec:pipeline}); the answer is found in
constant work after the initial Dijkstra, with the $6$-tuple enumeration
never firing.

At $k=15$ the ratio is
$56.18~\text{s} / 22~\text{\textmu s} \approx 2.6$~million$\times$ in
favour of CWZ. Past $k=15$ Yen is too slow for our budget; CWZ continues
to run unimpeded up to $k=24$ ($|V|=73$) in tens of microseconds.

\paragraph{A measurement note.} An earlier version of this benchmark
interleaved the two algorithms, timing CWZ at each $k$ immediately after
the Yen run at the same $k$. That inflated the CWZ timings at
$k=13,14,15$ by $15$--$50\times$ (to $0.4$--$1.4$~ms), because each CWZ
run began in the allocator and page-table state left behind by a Yen run
that had just churned $2^{k}$ paths. The symptom was a CWZ column that
rose with $k$ up to $k=15$ and then collapsed by a factor of $50$ at
$k=16$, where Yen no longer ran. The harness now measures the two
algorithms in separate passes, with a warm-up and $11$ repetitions per
CWZ point; the resulting CWZ column is flat-to-mildly-increasing across
the whole sweep, as the algorithm's structure predicts. The correction
moves the comparison further in CWZ's favour, but the original figures
were an artefact of the harness rather than of either algorithm.

\paragraph{Significance.} This is the empirical fact CWZ's worst-case
polynomial bound is buying. On any input where the simple heuristic
blows up, the polynomial-time algorithm continues to terminate quickly.
The \emph{benign-random} comparison in the previous sections must be
read in light of this: Yen is faster in practice on random inputs
because the adversarial structure does not arise by accident, but it
cannot be relied on as a worst-case solver. The diamond-chain panel of
Figure~\ref{fig:runtime} shows the same reversal in the routine
benchmark, side-by-side with the four benign families.

\section{Threats to validity}

\paragraph{Fidelity of the implementation.}
\begin{itemize}
  \item Our 2-VDP-in-DAG subroutine is a max-flow-based implementation
    rather than Tholey's linear-time
    algorithm~\cite{Tholey2012}. Both have the same asymptotic
    $O(|V|+|E|)$ bound, but the constant factor differs; we have not
    measured the gap.
  \item Our reduction to $(s,t)$-straight (\texttt{reduce\_to\_straight}
    in \texttt{src/cwz/reductions.cpp}) caps each multi-edge family at
    $K=5$ distinct weights per (source, destination) pair after each
    fold. This is the one place where our implementation is an
    empirically-validated approximation rather than a proven
    transcription: the paper's \textsc{NextSP} instead records
    candidates for expensive routes through a folded vertex and keeps
    only the minimum shortcut. Ablating $K$ against the brute-force oracle
    over $10{,}000$ random instances at $|V|\le 12$ ($7{,}138$ of which
    have an NSP) puts $K=1$ at $2{,}092$ misses ($29.3\%$) and
    $K=2$, $K=3$ and $K=5$ at none, so the margin we actually rely on is
    one extra weight, not five. But no
    constant $K$ is provable: a route can be unusable because the folded
    vertices it passes through are already used elsewhere on the path,
    and nothing bounds how many cheap alternatives such a collision can
    rule out. A pathological input whose only NSP needs the
    $6^{\text{th}}$ or later cheapest shortcut weight cannot be excluded
    a priori.
  \item The layered solver (\texttt{cwz\_nsp\_layered}) is a direct
    transcription of the paper's \textsc{NextSP-Layered}: the single
    $6$-tuple enumeration of Lemma~5.3, on a strictly $(s,t)$-layered
    graph, with the removed-back-edge candidates supplied by the
    reduction. This enumeration \emph{alone} reproduces the brute-force
    oracle on all $34{,}000$ trials.
  \item Both oracles are our own code. The brute-force search shares no
    modules with the CWZ pipeline -- it carries its own Dijkstra
    specifically so that a bug in \texttt{src/shortest\_path/} cannot
    hide inside the oracle that is supposed to catch it -- but it was
    written by one author from one reading of the problem statement. A
    shared misunderstanding of the NSP definition would be invisible to
    both. Yen, implemented independently from a published algorithm, is
    the partial mitigation, and it agrees on every benign instance.
\end{itemize}

\paragraph{Validity of the measurements.}
\begin{itemize}
  \item Runtimes are not reproducible to better than roughly a factor of
    three. The diamond-chain family ignores the random seed, so its
    $|V|=50$ instance is deterministic -- the same graph on every run --
    yet Yen measured $1.54$~s on one run of the sweep and $4.82$~s on
    another, the latter confirmed by five standalone repetitions on an
    idle machine. The main sweep takes a \emph{single} wall-clock
    measurement per instance, with no repetition, on a shared laptop
    under WSL2; only the adversarial sweep of
    Section~\ref{sec:adversarial} repeats each point. No runtime claim in
    this chapter finer than ``orders of magnitude'' should be relied on.
  \item Sub-microsecond runtimes fall below our timer resolution and are
    clamped to a $1$~\textmu s floor for display. The flat segments at
    that level in Figure~\ref{fig:runtime} mean ``too fast to measure'',
    not ``constant time'', and no comparison between algorithms is
    meaningful there.
  \item Twenty instances per cell is a moderate sample: a single instance
    moves a rate by $5$ percentage points. The small dip in
    Figure~\ref{fig:existence} at $|V|=30$ is one such instance and is
    noise -- at $200$ trials the same configuration gives $100\%$. The
    IQR bands on Figure~\ref{fig:runtime} expose the corresponding
    spread in runtime; on Erd\H{o}s--R\'enyi at large $|V|$ the band
    covers roughly a factor of two around the median.
\end{itemize}

\paragraph{Generality of the inputs.}
\begin{itemize}
  \item The sampled distribution is \emph{conditioned}. We redraw any
    instance in which $t$ is unreachable from $s$
    (Section~\ref{sec:methodology}), so what we measure is not
    Erd\H{o}s--R\'enyi $G_{n,p}$ but $G_{n,p}$ conditioned on $s\to t$
    connectivity. At small $|V|$ the two differ substantially -- at
    $|V|=6$ roughly two thirds of raw draws were being rejected. The
    conditioning is the right choice for measuring an NSP solver, since
    a graph with no $s\to t$ path poses no NSP question at all, but our
    results describe the conditioned distribution and do not transfer
    unchanged to the unconditioned one.
  \item The horizontal axis is the \emph{target} $|V|$, and the actual
    size differs by family: at target $50$ the grid instances have $64$
    vertices, the layered instances $51$, and the diamond chains $49$.
    Comparisons \emph{within} a family across sizes are sound;
    comparisons \emph{between} families at a fixed abscissa are not
    strictly like-for-like.
  \item The generator families we sweep -- excluding diamond-chain --
    are benign: none deliberately constructs many alternate shortest
    paths to stress Yen. The dedicated diamond-chain sweep
    (Section~\ref{sec:adversarial}) addresses this gap directly, but it
    is a single hand-built construction rather than a family of
    adversarial structures.
  \item Yen's enumeration cap $k_{\max}=5000$ is a configuration choice,
    and it is what fixes the point at which Yen stops merely being slow
    and starts returning the wrong answer (Section~\ref{sec:agreement}).
    That a capped Yen must eventually fail on graphs with exponentially
    many shortest paths is unavoidable; that it fails specifically at
    $|V|\ge 40$ on this family is an artefact of the cap we chose.
  \item All experiments stop at $|V|=50$ and use integer weights drawn
    uniformly from $[1,10]$. Narrow weight ranges make ties -- and hence
    equal-cost shortest paths -- more likely than a continuous or wider
    distribution would, which favours neither algorithm uniformly: ties
    are what defeat both Yen and the brute-force bound, while CWZ is
    indifferent to them.
\end{itemize}
